\documentclass[11pt,a4paper]{article}

\usepackage[T1]{fontenc}
\usepackage[utf8]{inputenc}
\usepackage{lmodern}
\usepackage[margin=3.1cm,top=3.0cm,bottom=3.2cm]{geometry}
\usepackage{amsmath,amssymb,amsthm,mathtools}
\usepackage{booktabs}
\usepackage{microtype}
\usepackage{fancyhdr}
\usepackage{enumitem}
\usepackage[hidelinks]{hyperref}

\linespread{1.08}
\setlength{\parindent}{1.4em}
\setlength{\parskip}{0pt}

\theoremstyle{plain}
\newtheorem{proposition}{Proposition}
\newtheorem{corollary}[proposition]{Corollary}
\theoremstyle{definition}
\newtheorem{definition}[proposition]{Definition}
\theoremstyle{remark}
\newtheorem{remark}[proposition]{Remark}

\pagestyle{fancy}\fancyhf{}
\renewcommand{\headrulewidth}{0.4pt}
\fancyhead[L]{\small\itshape The Liquidator's Frontier}
\fancyhead[R]{\small\thepage}

\newcommand{\LIF}{\mathrm{LIF}}
\DeclareMathOperator{\ltv}{LTV}

\begin{document}

\begin{titlepage}
\vspace*{7.5cm}
\begin{center}
{\LARGE\bfseries The Liquidator's Frontier}\\[0.9em]
{\large Leverage, Seniority and the Endogenous Attachment Point\\ in Isolated Lending Markets}\\[2.6em]
{\normalsize $\sigma$-Labs}\\[0.9em]
{\normalsize 10 September 2026}
\end{center}
\end{titlepage}

\newpage
\section*{Abstract}
\thispagestyle{fancy}

A collateralised borrow in an isolated lending market is not a loan with a haircut. It is a
two-tranche structure: the borrower holds the residual claim and is long a put on the collateral,
the lender holds a senior one. What is not standard is that the attachment point is of two kinds.
The \emph{contractual} attachment --- the price at which senior loss becomes arithmetically
unavoidable --- is fixed at creation and has a closed form. The \emph{effective} attachment ---
the price at which the mechanism stops working --- is produced block by block by an oracle, a
liquidation incentive and an interest-rate model, is never written down, and is the one that binds.

This paper derives the first in closed form for markets of the Morpho Blue type, then spends its
length on the second. Two corrections to the received account follow. The borrower's put is struck
at the senior threshold, not at the liquidation threshold, so a liquidated borrower is not a
wiped-out one; and the naive account, which ignores the premium paid to the agent who enforces the
structure, overstates both the equity cushion and the buffer by exactly that premium --- $2.9$
points of the entry price at production parameters, against a $7.9$-point interval between the two
thresholds.

The central result concerns not when liquidation is permitted but when it is performed. Modelling
a flash-financed liquidator, who receives collateral at the oracle price and must sell it at the
market price, gives one incentive inequality. It yields a maximum tolerable slippage $\sigma^{*}$
and a minimum viable position size proportional to gas. That budget is not free: $\sigma^{*}$ and
the room the liquidator has to work in are the same primitive $1-\lambda$, split in the fixed
proportion $c$ to $1-c$, so what pays the enforcing agent is taken out of the distance over which
it must succeed.

The inequality is then made dynamic. Slippage need not be estimated --- the venue state is readable
at every block --- so the frontier is a floor computable each block, free of any distributional
assumption, at the cost of being a best case. A position larger than one admissible slice must be
cleared in sequence, and we bound how many slices a market can carry by a number independent of the
borrower's leverage. Impact, volatility and carry then enter one first-passage problem whose
degenerate limits are the deterministic and driftless cases, and two markets with identical caps
and identical room are shown to differ by orders of magnitude in the probability of senior loss.
Finally we bound where a liquidation cascade stops --- with finite borrow it always stops --- by a
secant condition valid at the cliffs of a step function, where a contraction argument is not.

Three apparently distinct failure modes --- a reverting oracle, illiquid collateral, and dust
positions nobody bothers to close --- are the same inequality losing on different terms. We close
with the consequence for parameter choice: raising the loan-to-value cap does not only thin the
residual tranche, it shrinks the enforcing agent's slippage budget, so the cap is a claim about the
collateral's liquidity stated in the grammar of a risk parameter.

\medskip
\noindent\textbf{Keywords.} Collateralised lending; tranching; attachment point; liquidation
incentives; flash loans; market microstructure; decentralised finance; lending protocols.

\newpage
\section{The setting}

Fix an isolated lending market in the sense of Morpho Blue: one collateral asset, one loan
asset, one oracle, one interest-rate model, and one liquidation loan-to-value cap
$\lambda \in (0,1)$, all immutable at creation.

A borrower holds collateral worth $A$ in units of the loan asset, valued at the oracle price,
and owes $D$. Write
\[
  \ltv \;=\; \frac{D}{A}, \qquad E \;=\; A-D, \qquad L \;=\; \frac{A}{E} \;=\; \frac{1}{1-\ltv},
\]
so that $D = (L-1)E$ and $\ltv_0 = (L-1)/L$ at entry. The position is \emph{healthy} while
\begin{equation}\label{eq:healthy}
  A\lambda \;\geq\; D,
\end{equation}
which is the on-chain solvency predicate, evaluated with the oracle price and nothing else.

Throughout, $\ell$ denotes a proportional fall in the price of the collateral, so the collateral
is worth $A(1-\ell)$ after the move, and $D$ is unchanged over the horizon considered unless
interest is explicitly in play.

\subsection{Why this is a tranche}

Three positions on the same asset must be distinguished, because they are routinely conflated.

Holding the collateral outright gives a linear payoff with unbounded downside. Posting it and
borrowing gives something else: the borrower keeps the drawn $D$, and if the collateral falls
far enough it is seized and the debt is extinguished against it. The borrower's loss is
therefore truncated, and the lender's recovery is $\min(D,\,A/\LIF)$ for the incentive factor
$\LIF$ defined below --- which is $D$ less $1/\LIF$ puts on the collateral struck at $D\,\LIF$.
The lender is short those puts and the borrower is long them, which is what seniority means in
credit.

Two remarks on the strike, both of which the informal account gets wrong. It is $D\,\LIF$, the
level at which exhaustive seizure stops covering the debt --- not the liquidation threshold
$D/\lambda$, which is a barrier at which the mechanism switches on, not a level at which anyone
is paid. And the two are distinct: at the liquidation threshold, exhaustive seizure returns
$A - D\,\LIF > 0$ to the borrower. Being liquidated is not being wiped out.

Re-levering does not create the residual tranche; the first borrow already did. It thins it.

\section{Tranche geometry}

\begin{proposition}[Liquidation threshold]\label{prop:liq}
A position entered at $\ltv_0 \leq \lambda$ becomes liquidatable after a proportional price fall
\[
  \ell_{\mathrm{liq}} \;=\; 1-\frac{\ltv_0}{\lambda} \;=\; 1-\frac{L-1}{L\lambda}.
\]
\end{proposition}

\begin{proof}
By \eqref{eq:healthy} the position ceases to be healthy when $A(1-\ell)\lambda < D$, that is
$1-\ell < D/(A\lambda) = \ltv_0/\lambda$. Substituting $\ltv_0=(L-1)/L$ gives the second form.
\end{proof}

Enforcement is not free. A liquidator who repays $r$ receives collateral worth $r\,\LIF$, where
the incentive factor is fixed by two protocol constants.

\begin{definition}[Liquidation incentive factor]
With cursor $c$ and cap $M$,
\[
  \LIF(\lambda) \;=\; \min\!\left(M,\ \frac{1}{1-c\,(1-\lambda)}\right),
  \qquad c=0.3,\quad M=1.15 .
\]
\end{definition}

$\LIF$ is decreasing in $\lambda$: a tighter cap pays the enforcer more. The reason is a
constraint, not a judgement about risk --- the bonus has to fit inside the collateral surplus
$1/\lambda-1$ that still exists when liquidation opens, and that surplus shrinks as the cap rises.
Throughout we assume $M$ is not binding, which holds for $\lambda > 0.5652$ and therefore for every
parameter considered here.

\begin{proposition}[Senior threshold and loss]\label{prop:sen}
Exhaustive seizure of collateral worth $A$ repays at most $A/\LIF$. Consequently senior loss is
unavoidable if and only if $A < D\,\LIF$, that is after a fall
\[
  \ell_{\mathrm{sen}} \;=\; 1-\ltv_0\,\LIF ,
\]
and its magnitude is $D - A/\LIF$, not $D-A$.
\end{proposition}

\begin{proof}
A liquidator repaying $r$ removes $r\,\LIF$ of collateral value, so the largest repayment the
collateral can fund is $r_{\max}=A/\LIF$. If $A/\LIF \geq D$ the debt is fully repayable and
$A - D\,\LIF$ of collateral survives. Otherwise the residual debt $D-A/\LIF$ is unrecoverable.
Writing $A = A_0(1-\ell)$ and dividing by $A_0$ gives the threshold.
\end{proof}

The threshold $\ell_{\mathrm{sen}}$ is the \emph{contractual attachment point} of the structure: the
fall at which the senior claim begins to be impaired, determined by $\lambda$, fixed at creation,
and by $\ltv_0$, fixed when the position is opened and thereafter drifting only with carry. It is
available in closed form to anyone who reads the parameters. It is the object a
static tranching analysis would stop at. Sections \ref{sec:frontier} to \ref{sec:cascade} are
about the other one.

The correction matters at the boundary. At $A=D$ --- the point a naive account would call
break-even --- the lender has already lost $D(1-1/\LIF)$, which is $4.20\%$ of the debt at
$\lambda=0.86$ and $11.25\%$ at $\lambda=0.625$.

\begin{definition}[Buffer]
\[
  B \;=\; \ell_{\mathrm{sen}}-\ell_{\mathrm{liq}} \;=\; \ltv_0\left(\frac{1}{\lambda}-\LIF\right).
\]
\end{definition}

$B$ is the interval of collateral price over which liquidation is both permitted and sufficient:
the room the mechanism has to work in, measured from the entry price.

\begin{proposition}[The buffer is never empty]\label{prop:bpos}
$B>0$ for every $\lambda\in(0,1)$, every $\ltv_0>0$, every cursor $c<1$ and every cap $M$.
\end{proposition}

\begin{proof}
$B>0$ iff $\lambda\,\LIF<1$. Since $\LIF \leq 1/(1-c(1-\lambda))$ by definition of the minimum,
$\lambda\,\LIF \leq \lambda/\bigl((1-c)+c\lambda\bigr)$, and
$\lambda < (1-c)+c\lambda \iff (1-c)\lambda < 1-c \iff \lambda<1$. No case split on $M$ is needed.
\end{proof}

The positivity is a design constraint rather than a discovery: the incentive formula is built to
satisfy $\LIF < 1/\lambda$, which is what makes the structure a tranche at all rather than one in
which the enforcer is paid out of the lender's principal from the first block.

The buffer can be small. At $L=3$ and $\lambda=0.86$, liquidation opens at $-22.5\%$, senior loss
begins at $-30.4\%$, and $B=7.9$ points. The naive account --- which asks only when the equity is
extinguished, at a fall of $1/L=33.3\%$ --- is precisely the case $\LIF=1$. It places the
liquidation threshold correctly, since $\ell_{\mathrm{liq}}$ does not involve $\LIF$ at all, and
misses the senior threshold by
\[
  \tfrac{1}{L}-\ell_{\mathrm{sen}} \;=\; \ltv_0\,(\LIF-1),
\]
the enforcer's premium expressed as a fraction of the entry price --- $2.9$ points here. The
buffer inherits exactly that error, since $B$ is the difference of the two thresholds: the naive
counterpart of $B$ is $1/L-\ell_{\mathrm{liq}} = 10.8$ points against a true $7.9$. Note that
$1/L$ and $B$ are not themselves comparable quantities; the first is a cushion measured from
entry, the second an interval between two thresholds.

Two measures of the same room will both be needed below, and confusing them is the easiest error
to make here. $B$ is measured from the \emph{entry} price and therefore depends on $L$. Measured
from the price at which liquidation opens --- which is where the mechanism actually operates ---
the distance to the senior threshold is
\begin{equation}\label{eq:beta}
  \beta \;=\; 1-\lambda\,\LIF,
\end{equation}
which does not depend on $L$ at all. At $\lambda=0.86$, $\beta = 10.23\%$ while $B = 7.93\%$.
Every dynamic statement in this paper is about $\beta$.

\begin{proposition}[Partial liquidation below $\ell_{\mathrm{sen}}$]\label{prop:sign}
Let $r \in \bigl(0,\min(D,\,A/\LIF)\bigr)$. After a repayment $r$ the ratio becomes
$\ltv' = (D-r)/(A-r\,\LIF)$, and
\[
  \ltv' > \ltv \iff A < D\,\LIF .
\]
Moreover the eventual shortfall $D-A/\LIF$ is invariant under any such repayment.
\end{proposition}

\begin{proof}
$\ltv'>\ltv$ iff $(D-r)A > D(A-r\LIF)$, i.e. $-rA > -Dr\,\LIF$; dividing by $r>0$ and reversing
gives $A < D\,\LIF$. For the second claim,
$(D-r)-(A-r\LIF)/\LIF = D - r - A/\LIF + r = D-A/\LIF$.
\end{proof}

This is the structural content of Proposition~\ref{prop:sen}. Above the senior threshold each
partial liquidation lowers the loan-to-value ratio and the position converges to health; below
it, each partial liquidation \emph{raises} the ratio and the position never returns to health.
The threshold marks where the mechanism stops being able to repair a position, and it does so
before any bad debt is recorded.

It does not, however, mark a point at which liquidating destroys value, and the temptation to
read it that way should be resisted. The invariance in Proposition~\ref{prop:sign} says the
shortfall is fixed the moment the price crosses $\ell_{\mathrm{sen}}$: partial liquidation below
the threshold neither deepens nor cures it. If anything the incentive runs the other way. Every
unit repaid below the threshold converts an exposure that can still worsen into a loss that
cannot, so from the lender's side faster clearing is strictly better there --- the opposite of
what the rising ratio suggests. What changes at $\ell_{\mathrm{sen}}$ is only that the position
can now be worked to exhaustion rather than back to health.

\section{The liquidator's frontier}\label{sec:frontier}

Propositions \ref{prop:liq}--\ref{prop:sign} describe when liquidation is \emph{permitted}. A
lending market is cleared not by permission but by an agent who acts when acting pays.

Consider a liquidator financed by a flash loan, which is the standard construction and imposes no
capital requirement. The agent borrows $r$ of the loan asset, repays $r$ of debt, receives
collateral worth $r\LIF$ \emph{at the oracle price}, sells it into the market at proportional
slippage $\sigma$, returns the flash loan with fee $\phi$, and pays gas $g$:
\[
  \pi \;=\; r\,\LIF\,(1-\sigma) \;-\; r(1+\phi) \;-\; g .
\]

\begin{proposition}[Incentive inequality]\label{prop:frontier}
Liquidation of size $r$ is performed only if
\begin{equation}\label{eq:frontier}
  \LIF\,(1-\sigma) \;\geq\; 1+\phi+\frac{g}{r}.
\end{equation}
\end{proposition}

On Morpho Blue $\phi=0$: the flash-loan primitive transfers the assets out and pulls the same
amount back, charging nothing. We take $\phi=0$ below.

\begin{corollary}[Slippage tolerance]\label{cor:sigma}
In the limit of large $r$, the market clears only where
\[
  \sigma \;\leq\; \sigma^{*} \;=\; 1-\frac{1}{\LIF}.
\]
\end{corollary}

\begin{corollary}[Minimum viable size]\label{cor:rmin}
At gas $g$ and slippage $\sigma<\sigma^{*}$, no liquidation occurs below
\[
  r_{\min} \;=\; \frac{g}{\LIF\,(1-\sigma)-1},
  \qquad\text{so at }\sigma=0,\quad r_{\min}=\frac{g}{\LIF\,\sigma^{*}} .
\]
For $\sigma\geq\sigma^{*}$ no size is viable.
\end{corollary}

The second corollary is often quoted as a gas bound. It is not: $r_{\min}$ is proportional to gas
but inversely proportional to $\sigma^{*}$, and Table~1 shows it moving twentyfold across the
range of caps at fixed $g$. The cap sets the scale of the dust; gas only sets the unit.

\subsection{What the budget costs}

$\sigma^{*}$ is not a free parameter of the design. It is carved out of the same quantity that
supplies the buffer.

\begin{proposition}[Partition of the cap headroom]\label{prop:split}
Away from the cap $M$,
\[
  \sigma^{*} \;=\; c\,(1-\lambda), \qquad \beta \;=\; (1-c)(1-\lambda)\,\LIF,
\]
so that the headroom $1-\lambda$ splits as
\[
  1-\lambda \;=\; \underbrace{c\,(1-\lambda)}_{\text{enforcer's budget}} \;+\;
                  \underbrace{(1-c)(1-\lambda)}_{\text{room retained}},
  \qquad\text{and}\qquad
  \frac{\beta}{\sigma^{*}} \;=\; \frac{(1-c)\,\LIF}{c}.
\]
\end{proposition}

\begin{proof}
$\sigma^{*} = 1-1/\LIF = 1-\bigl(1-c(1-\lambda)\bigr) = c(1-\lambda)$. For $\beta$, use
\eqref{eq:beta}: $1-\lambda\LIF = \bigl(1-c(1-\lambda)-\lambda\bigr)\LIF = (1-c)(1-\lambda)\LIF$.
The ratio follows.
\end{proof}

This is the paper's central identity and it is worth stating without symbols. The cursor $c$ is a
\emph{split}. Every point of loan-to-value headroom the protocol declines to lend against is
divided between the premium that pays the enforcing agent and the room that agent is given to
work in, in the proportion $c$ to $1-c$. The retained room is a distance in loan-to-value;
converting it to a price distance multiplies by $\LIF$, which is why the ratio the liquidator
actually experiences, $\beta/\sigma^{*}$, is $(1-c)\LIF/c$ --- at production constants, roughly one
part budget to $2.4$ parts room. The two cannot be set independently: a protocol wishing to
pay liquidators more for a given cap can do so only by shortening the distance over which they
have to succeed. And since $\sigma^{*}=c(1-\lambda)$ is linear, the exchange rate is a constant:
\emph{every point of loan-to-value cap costs the liquidator thirty basis points of slippage
budget}, whatever the collateral and whatever the venue.

\begin{remark}[The budget is also the loss]
$\sigma^{*}$ is likewise the lender's proportional loss at $A=D$, since
$(D-A/\LIF)/D = 1-1/\LIF$ there. This is an identity rather than a theorem, but it is the reason
the two halves of this paper describe one object: the premium that pays the enforcer, the market
depth the collateral must have, and --- at the point $A=D$ --- the loss the lender absorbs when
that depth is absent are the same number.
\end{remark}

\begin{table}[h]\centering
\begin{tabular}{rrrrr}
\toprule
$\lambda$ & $\LIF$ & $\sigma^{*}$ & $\beta$ & $r_{\min}$ at $g=50$ \\
\midrule
0.980 & 1.0060 & 0.60\% & 1.41\% & 8{,}283 \\
0.965 & 1.0106 & 1.05\% & 2.48\% & 4{,}712 \\
0.945 & 1.0168 & 1.65\% & 3.91\% & 2{,}980 \\
0.915 & 1.0262 & 2.55\% & 6.11\% & 1{,}911 \\
0.860 & 1.0438 & 4.20\% & 10.23\% & 1{,}140 \\
0.770 & 1.0741 & 6.90\% & 17.29\% & 675 \\
0.625 & 1.1268 & 11.25\% & 29.58\% & 394 \\
\bottomrule
\end{tabular}
\caption{The frontier at production parameters. In loan-to-value units the split is the constant
$(1-c)/c = 2.33$; as a ratio of price distances $\beta/\sigma^{*}=(1-c)\LIF/c$ runs from $2.35$ to
$2.63$ across the range. $r_{\min}$ is in loan-asset units and
evaluated at $\sigma=0$; any positive slippage raises it.}
\end{table}

\begin{remark}[What $\sigma$ measures]
The collateral is transferred at the \emph{oracle} price and realised at the \emph{market} price,
so $\sigma$ is the oracle-to-execution gap and not merely order-book depth. A stale oracle
marking the collateral above the bid consumes the same budget as thin liquidity. An oracle that
reverts sets $\sigma=\infty$, since the trade cannot be initiated at all: the solvency predicate
\eqref{eq:healthy} cannot be evaluated, and the enforcement transaction fails before it reaches
the market. Withdrawal of collateral requires the same evaluation and freezes with it; lending is
not symmetrically frozen, as \S6 records.
\end{remark}

\section{Clearing}

Corollary~\ref{cor:sigma} bounds a rate. A lending market is exposed to a quantity. This section
turns one into the other.

\subsection{Slippage is observable, and part of its path is deterministic}

Fix the argument of $\sigma$ once, since it is easy to lose. Write $\sigma(v)$ for the
proportional slippage incurred selling collateral of loan-asset value $v$ into the venue, with
$\sigma$ increasing and $\sigma(0^{+})=0$. A liquidation repaying $r$ sells $r\,\LIF$, so the
frontier \eqref{eq:frontier} reads $\LIF\bigl(1-\sigma(r\LIF)\bigr)\geq 1+g/r$ throughout.

It is tempting to treat $\sigma$ as a random variable to be estimated. On an on-chain venue it is
not: the state is readable at every block, so $\sigma$ is a deterministic function of public
state at the moment of evaluation. Two qualifications are owed immediately. Liquidators
frequently exit through venues that are not on-chain, for which this is false; and the state at
the moment the frontier binds is not the state now, since liquidation follows an adverse move.

Part of that second gap is nonetheless deterministic. The debt is denominated in the loan asset
and does not fall with the collateral, so extinguishing it after a fall $\ell$ requires selling
$1/(1-\ell)$ times as many units. On a venue whose depth in units is unchanged this alone takes a
position from $2\%$ slippage to $3.3\%$ after a $40\%$ fall. The effect is real but its size is
model-dependent: on a constant-product venue the reserve of collateral grows as the price falls,
impact scales as $1/\sqrt{1-\ell}$ rather than $1/(1-\ell)$, and the same move gives $2.6\%$. The
honest statement is that the requirement grows with the move under any depth profile, and that
the growth rate has to be read off the venue rather than assumed.

\begin{remark}[The frontier gives a floor, computable each block]
A curator can evaluate \eqref{eq:frontier} at each block against a grid of hypothetical $\ell$
and obtain a threshold: the largest move this market absorbs while remaining enforceable. It
should be read as a best case, not a stress test. Current state plus the quantity effect is a
\emph{floor} on slippage --- it assumes the venue does not withdraw quotes --- and therefore an
\emph{upper bound} on the tolerable move. In stress, quote withdrawal routinely dominates the
quantity effect. The construction's value is that it is exact and requires no distributional
assumption; its limit is that the true threshold is always below it.
\end{remark}

\subsection{How many slices a market can carry}

\begin{definition}[Admissible slice]
The largest repayment a single profitable liquidation can carry, ignoring gas, is
\[
  q_{\max} \;=\; \sup\{\,q>0 \;:\; \sigma(q\,\LIF)\leq\sigma^{*}\,\},
\]
which is positive since $\sigma(0^{+})=0<\sigma^{*}$, and finite whenever $\sigma$ is unbounded.
\end{definition}

A position of size $Q>q_{\max}$ cannot be cleared in one transaction at a profit. It must be sold
in $N=\lceil Q/q_{\max}\rceil$ slices, separated by whatever interval $\Delta$ the venue needs for
arbitrage to restore the price. Write $T=N\Delta$ for the clearing time. Let $\rho\in[0,1]$ denote
the fraction of each slice's impact which is permanent, that is, not recovered between slices.

The liquidator therefore pushes the price toward the very threshold it is racing against. Two
opposing effects have to be tracked, and the first version of this argument tracked only one. Each
slice moves the price down by $\rho\sigma^{*}$, which consumes room. But each slice also repays
debt, and since seizure removes collateral value $\LIF$ times faster than it removes debt, the
\emph{proportional} distance from the current price to the senior threshold widens as the position
is worked. The bound below is stated conservatively, discarding the second effect.

Because impacts compound rather than add, the natural unit here is the log-decline. Write
\begin{equation}\label{eq:logunits}
  \kappa \;=\; -\ln\bigl(1-\rho\sigma^{*}\bigr), \qquad b \;=\; -\ln\bigl(\lambda\LIF\bigr),
\end{equation}
for the permanent log-impact of one slice and the log-distance from the liquidation price to the
senior threshold. To first order $\kappa \simeq \rho\sigma^{*}$ and $b \simeq \beta$; both are used
in this form for the rest of the paper.

\begin{proposition}[Slice-count bound]\label{prop:self}
Suppose clearing begins at the liquidation threshold, the price is unchanged except by the
liquidator's own permanent impact, and the oracle tracks that impact between slices. Then the
senior tranche survives the clearing of a position spanning $N$ slices whenever
\[
  N \;<\; \frac{b}{\kappa}, \qquad\text{to first order}\qquad
  N \;\leq\; N^{*} \;=\; \frac{\beta}{\rho\,\sigma^{*}} \;=\; \frac{(1-c)\,\LIF}{c\,\rho}.
\]
The bound is independent of the borrower's leverage and, up to $\LIF$, of the cap.
\end{proposition}

\begin{proof}
Let $x_n = A_n/\LIF$ be the repayment capacity of the collateral after $n$ slices and $D_n$ the
debt, so senior loss is the event $x_n < D_n$. A slice repaying $q$ seizes collateral worth
$q\LIF$ and then moves the price by a permanent factor $(1-u)$, $u=\rho\sigma^{*}$, so
$x_{n+1} = (x_n-q)(1-u)$ and $D_{n+1}=D_n-q$. Unrolling with equal slices $q=D_0/N$,
\[
  x_N \;=\; x_0(1-u)^N - q\sum_{k=1}^{N}(1-u)^k, \qquad D_N \;=\; D_0-Nq,
\]
so the slack $S_N = x_N-D_N$ satisfies
\[
  S_N \;=\; x_0(1-u)^N - D_0 + q\Bigl[N-\sum_{k=1}^{N}(1-u)^k\Bigr] \;\geq\; x_0(1-u)^N - D_0,
\]
the bracket being non-negative since $(1-u)^k\leq1$. Enforcement begins at $A_0=D_0/\lambda$, so
$x_0 = D_0/(\lambda\LIF)$ and $S_N \geq x_0\bigl((1-u)^N-\lambda\LIF\bigr)$, which is positive
exactly when $N\ln(1-u) > \ln(\lambda\LIF)$, that is $N<b/\kappa$. Replacing $1-(1-u)^N$ by $Nu$
gives the first-order form, and Proposition~\ref{prop:split} the closed form.
\end{proof}

The oracle hypothesis is not cosmetic and it cuts both ways. If the oracle tracks the impact, the
damage appears where the proof puts it: as movement of $A_n$ toward the senior threshold. If the
oracle lags, $A_n$ is marked too high, the position looks safer than it is, and the impact
surfaces instead as slippage on the next slice --- shrinking $q_{\max}$, raising $N$, and
tightening the same bound from the other side. A lagging oracle does not avoid the constraint; it
relocates it from the numerator to the denominator and conceals it while it accumulates.

At $\rho=0.3$ the bound is $b/\kappa=8.51$, or $N^{*}=8.12$ in first-order form: \emph{a position
more than about eight admissible slices in size cannot be assumed to survive its own unwinding,
even in a market that never moves.} The first-order constant is the partition ratio of
Proposition~\ref{prop:split} divided by $\rho$, which is why it does not depend on how levered the
borrower was --- a fact the entry-based buffer $B$ obscures.

The bound is sufficient, and slack. Under the same hypotheses, restoring the deleveraging effect
discarded in the proof --- the bracket dropped there --- and iterating to a full clear in equal
slices gives the sharp threshold: $17$ slices at $\rho=0.3$, $26$ at $\rho=0.2$. The gap is slack
in the bound, not uncertainty about the venue, and it decomposes cleanly: $8.12\to8.51$ is the
compounding restored by the logarithm, $8.51\to17$ is the deleveraging. The threshold does,
however, depend on how the position is worked, and in a direction worth stating. Under the proposition's own model --- impact
$\rho\sigma^{*}$ per slice whatever the slice's size --- splitting a slice in two doubles the
damage, and the sharp threshold falls from $17$ to $8$ at half-slices and $4$ at quarter-slices.
The coarsest admissible schedule is therefore optimal, and $17$ is the best case over schedules
rather than the only case. The conclusion survives a more realistic impact function so long as it
is concave or linear in size: with impact proportional to size the threshold runs from $17$ at
whole slices to $16.8$ in the continuous limit, so slicing moves it by about one per cent, and
with square-root impact finer slicing is again worse.
Only convex impact would let a liquidator do better by slicing, and that is not the usual shape.
What matters is that these numbers are small, and that they are properties of the venue rather
than of the loan.

This is the point at which the second attachment can be named.

\begin{definition}[Effective attachment point]\label{def:eff}
Write $d_{\mathrm{liq}}=-\ln(1-\ell_{\mathrm{liq}})$ and $d_{\mathrm{sen}}=-\ln(1-\ell_{\mathrm{sen}})$,
so that $d_{\mathrm{sen}}-d_{\mathrm{liq}}=b$ exactly. The \emph{effective attachment point} of a
position of size $Q$, with $N=\lceil Q/q_{\max}\rceil$, is
\[
  d_{\mathrm{eff}}(Q) \;=\; \max\bigl(d_{\mathrm{liq}},\; d_{\mathrm{sen}}-N\kappa\bigr),
  \qquad
  \ell_{\mathrm{eff}}(Q) \;=\; 1-e^{-d_{\mathrm{eff}}(Q)} ,
\]
a lower bound on the smallest exogenous fall from which clearing can no longer complete before the
senior claim is impaired.
\end{definition}

The definition is a rearrangement of the slack inequality in the proof above, not a new object:
the position must survive $N$ slices of permanent impact, so the fall it can absorb beforehand is
the contractual one less the damage the clearing will do. In log units that subtraction is exact,
and the guard against large $N$ is automatic --- $d_{\mathrm{sen}}-N\kappa$ simply falls through
$d_{\mathrm{liq}}$ and is clipped there.

Two properties follow, and neither was available for $\ell_{\mathrm{sen}}$. First,
$\ell_{\mathrm{liq}} \leq \ell_{\mathrm{eff}} < \ell_{\mathrm{sen}}$, the floor attained exactly
at $N=b/\kappa$ --- the two thresholds of \S2 are the two ends of this one quantity, and $b/\kappa$
is where the room has been entirely consumed by the act of clearing. The upper end is approached
but not attained, since even a single slice does $\kappa$ of damage: $\ell_{\mathrm{eff}}(q_{\max})
= 29.8\%$ against $\ell_{\mathrm{sen}}=30.4\%$ at the parameters below. That last $0.6$ points is
slack of the bound rather than a property of the position; the sharp recursion, which credits the
seizure before the impact, recovers it. Second, $\ell_{\mathrm{eff}}$ depends on $Q$, on $q_{\max}$
and on $\rho$, none of which is a parameter of the market and all of which move without any
transaction on the protocol; $\Delta$ enters too, once the price is allowed to move on its own,
through Proposition~\ref{prop:passage}. At $\lambda=0.86$, $L=3$ and $\rho=0.2$, a position of five
admissible slices has $\ell_{\mathrm{eff}}=27.4\%$ against a contractual $30.4\%$; at ten slices,
$24.3\%$.

This is the sense in which the attachment point of these structures is endogenous. It is not a
term of the contract but an outcome of the venue, recomputed every block, and it is the one that
binds. Substituting the sharp recursion for the slack bound gives
\[
  \ell_{\mathrm{eff}}^{\,\mathrm{sharp}}(N) \;=\;
  1-(1-\ell_{\mathrm{sen}})\,(1-u)\,\frac{(1-u)^{-N}-1}{Nu}, \qquad u=\rho\sigma^{*},
\]
which restores $\ell_{\mathrm{eff}}(q_{\max})=\ell_{\mathrm{sen}}$, halves the initial slope in $N$,
and moves the floor from $b/\kappa$ to $17$. What the two forms share is only the concave descent
between the two thresholds, once $N$ is measured as a fraction of its own threshold.

\subsection{Impact, volatility and carry in one race}

Proposition~\ref{prop:self} assumes the price moves only through execution. The opposite
idealisation --- that it moves only through news --- is equally available and equally partial.
Both are limits of one problem.

Write $Z_t = \ln\bigl(A_t/(D_t\LIF)\bigr)$. Senior loss is the event $Z_t \leq 0$, and
$Z_0 = \ln\bigl(1/(\lambda\LIF)\bigr) =: b$ at the moment enforcement begins. Three things move
$Z$: the liquidator's permanent impact, at rate $m=\kappa/\Delta$ with $\kappa$ from
\eqref{eq:logunits} --- $Z$ is a logarithm and one slice multiplies the price by $1-\rho\sigma^{*}$,
so $\kappa$ and not $\rho\sigma^{*}$ is the per-slice drop; the carry, at rate
$r_b-y$ where $r_b$ is the borrow rate and $y$ the collateral's yield; and exogenous price
innovation, with volatility $\sigma_p$. Repayment is omitted from $Z$, exactly as it is from the bound of
Proposition~\ref{prop:self}: by Proposition~\ref{prop:sign} each slice raises $Z$, so the
probability below is an upper bound, and its $\sigma_p\to0$ limit inherits the slack of $N^{*}$
rather than the sharp threshold of $17$.

\begin{proposition}[Probability of senior loss]\label{prop:passage}
Suppose $Z_t = b-\mu t+\sigma_p W_t$ with $\mu = m+(r_b-y)$, and let clearing take $T=N\Delta$.
Then
\[
  \mathbb{P}(\text{senior loss before } T)
  \;=\; \Phi\!\left(\frac{\mu T-b}{\sigma_p\sqrt{T}}\right)
      \;+\; e^{2\mu b/\sigma_p^{2}}\,
        \Phi\!\left(\frac{-\mu T-b}{\sigma_p\sqrt{T}}\right).
\]
\end{proposition}

\begin{proof}
Standard first passage for arithmetic Brownian motion: for $X_t=-\mu t+\sigma_p W_t$ and $a<0$,
$\mathbb{P}(\min_{t\leq T}X_t\leq a) = \Phi\bigl((a+\mu T)/(\sigma_p\sqrt T)\bigr) +
e^{-2\mu a/\sigma_p^{2}}\Phi\bigl((a-\mu T)/(\sigma_p\sqrt T)\bigr)$. Set $a=-b$.
\end{proof}

\begin{corollary}\label{cor:limits}
As $\sigma_p\to0$ the probability tends to $\mathbf{1}\{\mu T> b\}$, the value at $\mu T=b$ being
$\tfrac12$; ignoring carry, $\mu T =
N\kappa$ and the survival condition is $N<b/\kappa=8.51$, which is exactly
Proposition~\ref{prop:self}. As $\mu\to0$ it tends to
$2\Phi\bigl(-b/(\sigma_p\sqrt T)\bigr)$, the driftless reflection-principle bound.
\end{corollary}

Call $(b-\mu T)/(\sigma_p\sqrt T)$ the \emph{clearing margin}: room remaining after the
liquidator's own damage, divided by the distance the price is expected to wander while the
position is worked. Its content is that $\lambda$, and therefore $\beta$, is uninformative alone.
Two markets at the same cap, and therefore with identical buffers, can sit at opposite ends of the
outcome space.

\begin{table}[h]\centering
\begin{tabular}{lrr}
\toprule
& \textbf{deep, moderate vol.} & \textbf{thin, high vol.} \\
\midrule
$q_{\max}$ & 2{,}000{,}000 & 1{,}000{,}000 \\
position $Q$ & 10{,}000{,}000 & 10{,}000{,}000 \\
slices $N$ & 5 & 10 \\
$\Delta$ & 5 min & 15 min \\
$T$ & 25 min & 2.5 h \\
$\sigma_p$ & 100\%/yr & 300\%/yr \\
$\rho$ & 0.2 & 0.2 \\
\midrule
$b$ & 10.79\% & 10.79\% \\
impact $\mu T$ (carry omitted) & 4.22\% & 8.44\% \\
noise $\sigma_p\sqrt T$ & 0.69\% & 5.07\% \\
clearing margin & \textbf{9.53} & \textbf{0.46} \\
$\mathbb{P}(\text{senior loss})$ & ${\approx}10^{-21}$ & $41\%$ \\
\bottomrule
\end{tabular}
\caption{Identical cap, identical room $b$, opposite outcomes. Both positions satisfy
Proposition~\ref{prop:self} at $\rho=0.2$, so neither fails on impact alone. The difference is
carried entirely by $q_{\max}$, $\Delta$ and $\sigma_p$, none of which appears among the market's
parameters. The clearing margin, not the probability, is the result: the left-hand entry is a
Gaussian tail at nine and a half standard deviations and should not be read as a number. Under a
Student $t$ with three degrees of freedom the two margins give terminal tails of
${\approx}10^{-3}$ and $33\%$, against Gaussian terminal tails of $7\times10^{-22}$ and $32\%$ ---
a factor of a few hundred rather than of $10^{21}$, and the ordering unchanged. The $41\%$ above
includes the reflection term; the terminal tails do not.}
\end{table}

\begin{remark}[Where carry actually bites]
In Proposition~\ref{prop:passage} carry enters $\mu$ additively with impact, which places the
interest-rate model and the venue's depth in the same units --- but over a clearing window it is
negligible: at a $300\%$ spread, $2.5$ hours of carry is $9$ basis points against $840$ of impact.
Carry is not what loses the race. It is what starts it: at a flat price and $r_b>y$, the position
reaches the liquidation threshold at $t^{*}=\ln(\lambda/\ltv_0)/(r_b-y)$, and $r_b$ is set by an
adaptive model that raises it precisely when utilisation rises --- that is, when borrowers cannot
repay and lenders cannot withdraw. At $L=3$ a position with $t^{*}$ of over a decade at a $2\%$
spread has a $t^{*}$ of about a month at a $300\%$ one, with nothing about the position changed.
The interest-rate model does not decide whether liquidation succeeds; it decides how many
positions are queued at the frontier when it is tested.
\end{remark}

\section{Cascades}\label{sec:cascade}

Proposition~\ref{prop:self} treats one position pushing the price against itself. In a market with
many borrowers the same impact makes \emph{other} positions liquidatable, and the question is how
far the resulting sequence runs.

Declines must now be added across positions entered at different prices, so we work in logarithms
throughout this section, where impact compounds additively and no first-order approximation is
needed. Write $d=-\ln(1-\delta)$ for a log-decline measured from the pre-shock price, and recall
$\kappa$ and $b$ from \eqref{eq:logunits}. All positions are re-marked at the pre-shock price, so
$d_{\mathrm{liq}}(i)$ is position $i$'s liquidation threshold expressed in that common frame.

Let $d_0>0$ be an exogenous log-decline and let $Q(d)$ be the borrow rendered liquidatable once
the price has fallen by $d$. $Q$ is a primitive of the market, read off the distribution of
liquidation prices across open positions; it is a non-decreasing step function, with a jump
wherever a mass of positions shares a liquidation price. Clearing that volume obeys the mechanics
of \S4.2 --- $Q(d)/q_{\max}$ slices, each leaving $\kappa$ of permanent log-impact --- so
\begin{equation}\label{eq:cascade}
  d_{n+1} \;=\; F(d_n), \qquad
  F(d) \;=\; d_0+\kappa\,\frac{Q(d)}{q_{\max}},
  \qquad d_1 = F(d_0).
\end{equation}

\begin{proposition}[Where a cascade stops]\label{prop:cascade}
A least $d\geq d_0$ satisfying
\begin{equation}\label{eq:secant}
  Q(d) \;\leq\; \frac{q_{\max}}{\kappa}\,\bigl(d-d_0\bigr)
\end{equation}
exists whenever total borrow is finite, with no further hypothesis on $Q$. If in addition the
market holds finitely many positions, say $m$, the iterates \eqref{eq:cascade} reach it in at most
$m+1$ rounds, and the set of positions liquidated grows strictly at every round but the last.
\end{proposition}

\begin{proof}
$Q$ is non-decreasing, so $F$ is non-decreasing and maps the interval
$I=[d_0,\,d_0+\kappa Q_{\mathrm{tot}}/q_{\max}]$ into itself, $Q_{\mathrm{tot}}$ being total borrow.
$I$ is a complete lattice and \eqref{eq:secant} is $F(d)\leq d$ rearranged, so by Tarski's fixed
point theorem the set of fixed points of $F$ in $I$ is a non-empty complete lattice and has a least
element, which is the least $d$ satisfying \eqref{eq:secant}. Neither continuity nor finiteness is
used.

For the second claim, let $S_k=\{i : d_{\mathrm{liq}}(i)\leq d_k\}$ be the set liquidated after $k$
rounds. The iterates are non-decreasing --- $d_1\geq d_0$ since $Q\geq0$, and monotonicity of $F$
propagates it --- so the $S_k$ are nested increasing. If $S_{k+1}=S_k$ then $Q(d_{k+1})=Q(d_k)$,
since $Q$ depends on $d$ only through that set, hence $d_{k+2}=d_{k+1}$ and the iteration has
stopped. So every round but the last adds at least one position, and there are $m$ to add.
\end{proof}

The two halves of the proposition need different hypotheses, and separating them is the point.
Existence is a lattice fact and costs nothing. What finiteness buys is that the \emph{iteration}
reaches the fixed point, and reaches it fast --- a market of a thousand positions settles in at
most a thousand and one rounds, and in simulation far fewer. Without finiteness the limit of the
iterates need not satisfy
\eqref{eq:secant}: a right-continuous $Q$ with infinitely many jumps accumulating at some $d^{*}$
lets the sequence converge to $d^{*}$ while $F(d^{*})>d^{*}$, so the least $d$ satisfying
\eqref{eq:secant} lies strictly above the limit. Real markets have finitely many positions and the
question does not arise; a continuum model of the same market would need $Q$ left-continuous
instead.

\begin{remark}[Relation to network clearing]\label{rem:en}
Proposition~\ref{prop:cascade} is, in method, the clearing problem of Eisenberg and
Noe~\cite{eisenberg2001}. Their clearing vector solves $p=\bar L \wedge (e+\Pi^{\!\top}p)$, with
$\bar L$ the vector of total liabilities, $e$ the endowments and $\Pi$ the relative liability
matrix. The map is monotone on a lattice, so its solutions form a complete lattice and existence
is a lattice fact --- Tarski~\cite{tarski1955}, not Banach --- and their fictitious default
algorithm terminates in at most one round per node, because each round either adds a defaulter or
stops. Both arguments are the ones used above, and the second is why
the bound is $m+1$ rather than merely ``eventually''.

The economics is not theirs, and the difference matters. In~\cite{eisenberg2001} node $i$ reaches
node $j$ through a \emph{contractual liability}: $\Pi$ is a balance sheet. Borrowers in an isolated
market have no claim on one another whatsoever; position $i$ reaches position $j$ only by moving a
price, which is fire-sale contagion, and the closer ancestor is Cifuentes, Ferrucci and
Shin~\cite{cifuentes2005}, who layer a price that falls with aggregate sales onto the Eisenberg--Noe
problem and obtain a joint fixed point in payments and price. The reduced-form version used by
Greenwood, Landier and Thesmar~\cite{greenwood2015} and the estimates in Duarte and
Eisenbach~\cite{duarte2021} are the empirical descendants.

Two further departures are worth naming. Their state is a vector on a general liability graph;
ours collapses to a scalar, because in an isolated market every position reaches every other
through exactly one price --- the network is a star through the venue, not a graph, and that is a
property of the market design rather than a modelling convenience. And their algorithm iterates
\emph{downward} from full payment and selects the greatest clearing vector, where
\eqref{eq:cascade} iterates \emph{upward} from the shock and selects the least fixed point above
$d_0$. The direction is not a choice: a cascade genuinely begins at the shock and works up, so the
least fixed point is the one the market visits.
\end{remark}

Two things follow from the form of \eqref{eq:secant}. It is a secant condition, not a derivative
one, so it needs no smoothness and is evaluated exactly at the cliffs of $Q$, which are the only
places the question is interesting; a contraction argument would be both unavailable and
unnecessary. And it is a statement about \emph{where} the cascade stops, not whether. In a market
with finite borrow it always stops. The useful question is never termination --- it is whether
$\bar d$ arrives before the senior thresholds of the positions in its path.

Read \eqref{eq:secant} against Proposition~\ref{prop:self} and the two agree. A single position
starting at its own liquidation threshold makes $\bar d-d_0=N\kappa$, and its senior claim is
impaired only if that exceeds $b$ --- that is, only if $N>b/\kappa=8.51$, which is
Proposition~\ref{prop:self}'s $N^{*}=8.12$ up to the compounding the logarithm now carries
exactly. The market case is the same inequality with $Q$ read off many borrowers instead of one.
The quantity $q_{\max}/\kappa$ is the venue's \emph{marginal depth} in this model: the borrow that
may be made liquidatable per unit of log-decline without the decline feeding itself. The condition
reads

\begin{center}
\emph{the borrow made liquidatable over an interval of decline must not exceed the venue's
marginal depth times the width of that interval.}
\end{center}

Both sides are observable. The left is a property of how borrowers are distributed; the right is
a property of the market they will be sold into.

Two consequences deserve statement. First, the condition is \emph{local}: it can hold on average
across the price range and fail on an interval. A market in which many positions share a
liquidation price has a jump in $Q(\cdot)$ there, and a cascade that reaches that jump is carried
past it regardless of how comfortably \eqref{eq:secant} holds elsewhere. It is only reached if no
earlier $d$ satisfies \eqref{eq:secant} first, which is precisely why the criterion must be
evaluated as a secant from $d_0$ rather than pointwise. Concentration of \emph{liquidation
prices}, not merely of size, is the hazard --- and a front end offering one default leverage
produces it whenever entries also cluster in time, since a common leverage and a common entry
price give a common liquidation price.

Second, the condition couples two things that no single agent controls. A curator sets caps and
therefore influences $Q$; nobody sets $\sigma$. A market can pass the test on the day it is
configured and fail it after liquidity migrates, with no transaction on the protocol having
occurred.

Finally, \eqref{eq:secant} delivers the effective attachment point at market scale. Borrower $i$
can suffer senior loss only if the terminal $\bar d$ exceeds $d_{\mathrm{liq}}(i)+b$, which in
logarithms is exactly $d_{\mathrm{sen}}(i)$, with no first-order term discarded. The condition is
necessary and not sufficient: a position cleared early in the cascade is sold before the price
reaches its own senior level. Since $\bar d$ is non-decreasing in $d_0$, the market-scale effective
attachment of position $i$ is well defined as the smallest shock that carries the cascade past it,
$d_{\mathrm{eff}}(i)=\inf\{d_0 : \bar d(d_0) > d_{\mathrm{sen}}(i)\}$, with
$\ell_{\mathrm{eff}}(i)=1-e^{-d_{\mathrm{eff}}(i)}$ in the units of \S2. When $Q$ has a single step
of height $Q$ this is $d_{\mathrm{sen}}-\kappa Q/q_{\max}$, which is Definition~\ref{def:eff}
exactly. The contractual attachment of
Proposition~\ref{prop:sen} is a property of that position's parameters alone; the effective one is
a property of every other borrower
in the same market and of the venue they will all be sold into. Two positions on identical terms,
in two markets with identical parameters, do not have the same attachment point --- the one whose
neighbours' liquidation prices are more concentrated has the worse one, and nothing either
borrower signed records the difference.

\section{One inequality, three failures}

The three ways a market of this type stops behaving as a tranche structure are not independent
phenomena. They are \eqref{eq:frontier} failing --- on two terms, one of which fails in two
qualitatively different ways.

\begin{table}[h]\centering
\begin{tabular}{@{}l l p{5.4cm}@{}}
\toprule
\textbf{term} & \textbf{failure} & \textbf{observed as} \\
\midrule
$\sigma$ undefined & no price exists & liquidation and collateral withdrawal both impossible;
borrowing halted; only un-borrowed liquidity remains withdrawable \\
$\sigma>\sigma^{*}$ & collateral unsellable & bad debt on gapping or illiquid assets \\
$A/\LIF<r_{\min}$ & position too small & dust nobody clears; loss recognised late or never \\
\bottomrule
\end{tabular}
\end{table}

The first row is not a limiting case of the second. An oracle that reverts does not make the
collateral expensive to sell; it makes the solvency predicate \eqref{eq:healthy} unevaluable, so
the transaction fails before reaching any venue. Note also what it does \emph{not} freeze: lenders
whose funds are not currently borrowed can still withdraw, and typically do, which raises
utilisation on the remainder and --- by the interest-rate model --- the rate charged to borrowers
who cannot withdraw collateral. The freeze is asymmetric.

The third row is stated on collateral rather than on debt, since a liquidator can never repay more
than $A/\LIF$; when even that maximum repayment is below $r_{\min}$ no profitable liquidation
exists at any size, and ``never cleared'' follows without hedging. It deserves emphasis because it
inverts a common intuition. Large positions are not the dangerous ones \emph{on this term}: they
are the ones the gas condition is certain to clear --- though, by Proposition~\ref{prop:self},
they may well fail the other. The two size conditions run in opposite directions, $r_{\min}$ from
below and $N^{*}q_{\max}$ from above, so a market is enforceable only on the interval between
them, and a sufficiently thin venue closes that interval entirely.

A protocol that records loss only at exhaustive seizure --- as Morpho Blue does, writing bad debt
when and only when the remaining collateral reaches zero --- therefore does not recognise the loss
on precisely the positions that fail this term. The accounting is correct at every step and lags
reality for as long as $A/\LIF<r_{\min}$ holds, which, absent intervention, is forever.

\section{What the cap asserts}

The loan-to-value cap is read as a leverage limit. Proposition~\ref{prop:split} gives it a second
reading, and unlike the first it is quantitative.

Since $\sigma^{*}=c(1-\lambda)$, raising the cap thins the residual tranche and narrows the
enforcing agent's slippage budget in strict proportion, at thirty basis points of budget per point
of cap. At $\lambda=0.98$ the market clears only if the collateral can be sold within sixty basis
points of its oracle price; at $\lambda=0.86$ the budget is $4.2$ points. Neither number is a
statement about volatility, and neither is enforced anywhere in the protocol --- $\sigma^{*}$
appears in no contract, and nothing reverts when it is exceeded. What happens instead is that
nobody sends the transaction.

Setting a cap at $\lambda$ is therefore an assertion --- untested when made, and tested only in the
worst state --- that the collateral will trade within $c(1-\lambda)$ of its oracle price at the
size the market has accumulated. Read that way, a high cap is not a permission granted to
borrowers but a claim made on their behalf about a venue, and a low cap is better read as a
disclosure about the collateral than as a control on the borrower.

\section{What this does not settle}

The model treats $\sigma$ as a primitive and the collateral price as exogenous to execution beyond
the impact terms. Four things are parametrised rather than derived, and each is a distinct piece of
work.

\emph{Impact persistence and recovery.} Propositions~\ref{prop:self} and \ref{prop:passage} carry
$\rho$ and $\Delta$ as constants. Both are properties of the venue and of the information content
of the flow --- a liquidation is uninformed selling, which should recover faster than informed ---
and both are estimable per venue. Deriving them requires a market microstructure model this paper
does not supply, and $N^{*}$ is only as good as $\rho$.

\emph{The shape of $\sigma$ on a concentrated-liquidity venue.} The results use only that $\sigma$
is increasing; the numerical illustrations would change with the depth profile. On a venue where
liquidity is placed in ranges, $\sigma$ is not smooth and $q_{\max}$ can move discontinuously as
the price crosses a range boundary --- which is precisely the region where the frontier is tested.
This deserves separate treatment.

\emph{Competition among liquidators.} $\pi>0$ is necessary for one agent to act; the realised
clearing rate is set by the fastest, and priority auctions redistribute the premium without moving
the frontier. Whether competition shortens $\Delta$ enough to matter is an empirical question about
block construction rather than about lending.

\emph{Endogeneity of $r_b$.} The rate responds to utilisation, and utilisation responds to whether
the market clears. Proposition~\ref{prop:passage} takes $r_b$ as given; the fixed point of rate,
utilisation and exit behaviour is the natural object and remains open.

None of these weakens the frontier itself. \eqref{eq:frontier} is a necessary condition, so a market
failing it fails to clear regardless of what the extensions add. The conditions derived from it run
the other way: Propositions~\ref{prop:self} and \ref{prop:cascade} give sufficient conditions for
survival and for the location of the stopping point, and a market satisfying them may still fail
on a term this model holds fixed.

\begin{thebibliography}{9}
\bibitem{cifuentes2005} R.~Cifuentes, G.~Ferrucci and H.~S. Shin.
  Liquidity risk and contagion. \emph{Journal of the European Economic Association}, 3(2--3):556--566, 2005.
\bibitem{duarte2021} F.~Duarte and T.~M. Eisenbach.
  Fire-sale spillovers and systemic risk. \emph{Journal of Finance}, 76(3):1251--1294, 2021.
\bibitem{eisenberg2001} L.~Eisenberg and T.~H. Noe.
  Systemic risk in financial systems. \emph{Management Science}, 47(2):236--249, 2001.
\bibitem{greenwood2015} R.~Greenwood, A.~Landier and D.~Thesmar.
  Vulnerable banks. \emph{Journal of Financial Economics}, 115(3):471--485, 2015.
\bibitem{tarski1955} A.~Tarski.
  A lattice-theoretical fixpoint theorem and its applications.
  \emph{Pacific Journal of Mathematics}, 5(2):285--309, 1955.
\end{thebibliography}

\vspace{2em}
\hrule height 0.4pt
\vspace{0.8em}
{\footnotesize\noindent $\sigma$-Labs \ $\cdot$ \ \href{https://www.sigmalabs.fi}{sigmalabs.fi}}

\end{document}
